\section{Discussion}
\label{sec:discussion}

\subsection{Why correlation outperforms IG here}
Three explanations fit the ranking result.
\begin{enumerate}
  \item \textbf{Different questions.} IG asks which bins would change a few
        sampled pixels' component probabilities. Correlation asks which ion
        images follow a component across the whole section, which is closer to
        what mSCF1 rewards.
  \item \textbf{Sampling.} IG averages over 12 sampled pixels per component,
        while correlation uses every pixel. More pixels improved IG only
        slightly ($\le +0.024$) on the two sections tested, so sampling
        explains little of a gap of 0.05--0.25, but other IG settings
        (baselines, path counts, signed attributions) remain untested.
  \item \textbf{Metric alignment.} A segmentation that resembles the masks
        produces posterior maps that a mask-based metric naturally rewards.
\end{enumerate}
Correlation is therefore best read as an ablation finding about this pipeline:
the label-free segmentation carries the useful signal, and gradient attribution
through a dense VAE does not add to it. We do not claim that IG is generally
inferior, nor that segment-then-correlate, which is established
practice~\citep{alexandrov2011efficient}, is a new method.

\subsection{What remains of the spatial hypothesis}
The original hypothesis was that neighbourhood input would improve a VAE's
representation and therefore its peaks. In a dense VAE that averages or attends
over neighbours, no intervention gave a consistent gain beyond what nonlocal
(shuffled) context also provides, and smoothing the segmentation did not help
either. The attention experiments explain why. At $3\times3$ scale, adjacent
spectra are near-copies of the centre, so there is little to select between.
Spatial information did help in two other places: as a ranking statistic
(Moran's~$I$, and correlation with spatially coherent segments), and in S3PL,
which collapses without neighbour input. Even there, though, correctly placed
neighbours added only an inconsistent $+0.04$ over shuffled ones. Spatially
aware \emph{ranking}, rather than spatially aware \emph{encoding}, is what the
evidence supports, and it gives this report its title.

\subsection{Reconstruction is not ranking quality}
PCA-5 reconstructed GBM better than the VAE yet ranked peaks about equally
well, while the VAE segmented slightly better. Earlier development runs showed
the same pattern: context models reconstructed slightly worse but changed which
peaks were chosen, and a spatial penalty smoothed the latents while worsening
reconstruction. Reconstruction error rewards fitting all intensity variation,
including intense ions unrelated to tissue classes, whereas peak quality here
depends on how well a low-dimensional segmentation separates the classes. This
agrees with \citet{mazarura2025autoencoder}.

\subsection{S3PL instability and the reproduction gap}
The S3PL seed analysis shows that comparisons with a single S3PL run can be
misleading. Specific sections switch between distinct solutions depending on
the initial weights. The released preprocessing appears to create this
instability, since TIC normalisation removes it while also removing the
highest-scoring outcomes. Neither seeds nor normalisation explain the
published GBM score. We therefore report the gap as unresolved, and we do not
select favourable seeds or checkpoints for any baseline.

\subsection{Autonomy and the proposal's aim}
The proposal aimed for autonomous peak detection that removes human bias from
preprocessing, with a lightweight model as an efficient alternative to 3D
convolution. That aim is met only in part.
\begin{itemize}
  \item The pipeline fits and ranks without labels.
  \item Two settings still come from annotation: the number of segments and the
        peak budget. BIC could not replace the first.
  \item On cost, the 1D model used comparable time and less memory than S3PL
        on the section measured. It did not, however, gain from its spatial
        input, which was the reason for building it.
\end{itemize}

\subsection{Limitations and threats to validity}
\begin{itemize}
  \item \textbf{Statistical units.} GBM has four PDX lines, CAC one biopsy and
        RCC eight subjects. P-values describe these sections, not a population.
        With $n=8$, only small families can reach Holm significance, so many
        comparisons rest on practical labels and effect sizes.
  \item \textbf{Class-informed configuration.} The number of GMM components
        equals the number of annotated classes and was not supported by BIC.
        Peak budgets come from published, annotation-related counts. The
        pipeline is label-free in fitting and ranking only.
  \item \textbf{The metric.} mSCF1 measures agreement with expert masks, not
        molecular identity. Methods that produce mask-like maps are favoured.
  \item \textbf{RCC.} The three-class result partly reflects separating
        annotated tissue from a large unannotated region.
  \item \textbf{Post hoc elements.} Tests~S, V, N, R and Investigation~21 were
        chosen after earlier results were known. Each was predeclared before its
        own outputs, and all are reported, including null and negative
        outcomes. Recommending correlation is a post hoc conclusion.
  \item \textbf{Scope.} The VAE uses fixed-window, dense neighbourhood input.
        Graph-based and ion-image-level spatial models, cross-section transfer,
        and other attribution settings were not tested.
\end{itemize}
